\appendix

\section{Finite-size gap analysis}
\label{sec:gap}

We take the gap in the $P=-2$ symmetry sector,
\begin{equation}
\Delta_{\rm sec}(L,s)=E_1^{(P=-2)}(L,s)-E_0^{(P=-2)}(L,s),
\end{equation}
where $E_0^{(P=-2)}$ ($E_1^{(P=-2)}$) is the ground-state (first-excited-state)
energy in that sector.
The sector gap at $\delta=-2$ for $L=8,12,16,20,24$ is shown in
Fig.~\ref{fig:D02a}(a). Panel~(b) zooms to $s\in[0.53,0.65]$
(Fig.~\ref{fig:D02b}), marking the minimum of each curve: as $L$
grows the minima move toward $s=0.5$ and their values keep
decreasing, consistent with convergence to $s=0.5$ in the
thermodynamic limit. The gap at $s=0.5$ versus $1/L$
is shown in Fig.~\ref{fig:D02c}(c); the least-squares quadratic fit in $1/L$ over the five sizes has intercept $a=-0.017$ in units of the exchange coupling, small against the gap scale and consistent with
the gap-closing trend. These results characterize the finite-size gap evolution near $s=0.5$ at $\delta=-2$.
Ground states throughout are obtained by exact diagonalization
(Lanczos), and the minima analysis follows standard finite-size
practice for quantum spin systems, for which we follow the lecture
notes of Ref.~\cite{sandvik2010computational}.

\begin{figure*}[htbp]
\centering
\includegraphics[width=\textwidth]{{../../figs/exp03/exp03_D02}}
\caption{Finite-size gap analysis at $\delta=-2$: (a) symmetry-sector gap
$\Delta_{\rm sec}$ for $L=8,12,16,20,24$; (b) zoom to
$s\in[0.53,0.65]$ with curve minima marked, moving toward $s=0.5$ as
$L$ grows; (c) gap at $s=0.5$ versus $1/L$ for five
sizes with the least-squares quadratic fit in $1/L$; the square
marks the fitted intercept $a=-0.017$.
\label{fig:D02a}\label{fig:D02b}\label{fig:D02c}}
\end{figure*}

\section{Premium-qubit screening}
\label{sec:premium}

To run on hardware we screen candidate chains (rings) with two
quantities: the stabilizer mean, probing two-qubit gate quality, and
the readout fidelity, probing measurement quality.
The ideal output is the cluster state
$|C\rangle=\prod_{\langle ij\rangle}CZ_{ij}|+\rangle^{\otimes n}$,
where $\langle ij\rangle$ runs over the edges of the candidate subgraph;
its stabilizers $S_i$ (two-body $X_0Z_1$, $Z_{n-2}X_{n-1}$ at open ends,
three-body $Z_{i-1}X_iZ_{i+1}$ otherwise) have expectation $+1$ in the
absence of errors, so deviations of $\langle S_i\rangle$ from $1$ characterize combined state-preparation and measurement errors.
Each candidate is characterized by $4$ benchmark circuits: two
stabilizer circuits measuring the even- and odd-index stabilizers
separately (the stabilizers commute, while the two groups use different local measurement bases),
and two readout circuits without $CZ$ gates yielding the readout
fidelity $F_{\rm ro}$ from the mean of $P(00\cdots0)$ and
$P(11\cdots11)$. The score
\begin{equation}
{\rm score}=0.8\,\bar{S}+0.2\,F_{\rm ro},
\end{equation}
where $\bar{S}=n^{-1}\sum_i\langle S_i\rangle$ is the stabilizer mean,
ranks the candidates, weighting two-qubit gate quality over readout.
Cluster states~\cite{raussendorf2001oneway}
are the ideal outputs whose stabilizers we measure.
The $8$-chain and $10$-ring benchmarks differ (the ring includes a
closing $CZ$ and has three-body stabilizers throughout), so candidate rankings are computed separately within each geometry. The benchmark evaluates candidate subgraphs through cluster-state preparation and local measurements on the specified connectivity. The resulting rankings and selected qubits are reported below.

\section{Premium-qubit rankings}
\label{sec:rankings}

The screening on the Baihua processor yields three rankings. In the
$8$-chain geometry the top ten candidates are headed by the
$8$-qubit chain on qubits 77 through 70 with score $0.7801$ (one further
candidate tied within $0.02$). In the $10$-ring geometry the top ten are
headed by the $10$-qubit ring on qubits 69--73 and 82--86 with score
$0.5246$ (two further candidates tied within tolerance). Among the ten
$8$-subchains of the top-ranked ring, the subchain on qubits
69--73 and 82--84 is first, with score $0.7108$ taken from the
chain-benchmark measurements. Subsequent hardware experiments use the
leading candidate of each ranking.


